\documentclass[11pt]{article}

\usepackage{amsmath, amssymb, amsthm}
\usepackage[margin=1in]{geometry}

\DeclareMathOperator{\Var}{Var}
\DeclareMathOperator{\Bernoulli}{Bernoulli}

\title{Exponential Inequalities and Concentration of Measure}
\author{DeepSeek prompted by Riyad Derguini}
\date{}

\begin{document}

\maketitle

\textbf{Notation and Conventions.} Let \((\Omega, \mathcal{F}, \mathbb{P})\) be a probability space. For a random variable \(X\), we denote its expectation by \(\mathbb{E}[X]\) and its variance by \(\Var(X) = \mathbb{E}[X^2] - (\mathbb{E}[X])^2\). The Moment Generating Function (MGF) of \(X\) is \(M_X(t) = \mathbb{E}[e^{tX}]\), defined for \(t\) in a neighborhood of zero. The natural logarithm is denoted by \(\log\). We write \(X \sim \mathcal{N}(\mu, \sigma^2)\) to indicate that \(X\) follows a normal distribution with mean \(\mu\) and variance \(\sigma^2\).

\hrulefill

\section{Fundamental Scalar Inequalities}

The following analytic inequalities serve as the primitive tools from which probabilistic bounds are derived.

\subsection{The Exponential-Cosh Inequality}
For all \(x \in \mathbb{R}\):
\[
 \frac{e^x + e^{-x}}{2} = \cosh(x) \le e^{x^2/2}. 
\]

\emph{Proof.} Consider the Maclaurin series of both functions:
\[
 \cosh(x) = \sum_{k=0}^{\infty} \frac{x^{2k}}{(2k)!}, \qquad e^{x^2/2} = \sum_{k=0}^{\infty} \frac{x^{2k}}{2^k k!}. 
\]
It suffices to show that for all \(k \ge 0\), \((2k)! \ge 2^k k!\). For \(k=0\), equality holds (\(1=1\)). For \(k \ge 1\), observe that \((2k)! = (1 \cdot 2 \cdot \dots \cdot k) \cdot ((k+1) \dots (2k)) = k! \prod_{j=1}^k (k+j)\). Since \(k+j \ge 2\) for all \(j \ge 1\), we have \(\prod_{j=1}^k (k+j) \ge 2^k\). Thus, \((2k)! \ge 2^k k!\), and the inequality follows term-by-term. \(\square\)

\subsection{Basic Exponential Bounds}
For \(x \in \mathbb{R}\):
\begin{enumerate}
    \item \(1 + x \le e^x\) (with equality iff \(x=0\)).
    \item \(\log(1+x) \le x\) for \(x > -1\).
    \item \(e^x \ge 1 + x + \frac{x^2}{2}\) for \(x \ge 0\).
\end{enumerate}

\hrulefill

\section{Moment Generating Functions and the Chernoff Method}

The Chernoff method (or Chernoff-Cramér bound) is the standard technique for bounding tail probabilities using the MGF.

\subsection{The Chernoff Bound}
Let \(X\) be a random variable and \(a \in \mathbb{R}\). For any \(t > 0\),
\[
 \mathbb{P}(X \ge a) = \mathbb{P}(e^{tX} \ge e^{ta}) \le \frac{\mathbb{E}[e^{tX}]}{e^{ta}} = \frac{M_X(t)}{e^{ta}}. 
\]
Similarly, for \(t < 0\), \(\mathbb{P}(X \le a) \le \frac{M_X(t)}{e^{ta}}\).

\subsection{MGFs of Standard Distributions}
\begin{itemize}
    \item \textbf{Rademacher} (\(X \in \{-1, 1\}\) equiprobably): \(M_X(t) = \frac{e^t + e^{-t}}{2} = \cosh(t) \le e^{t^2/2}\).
    \item \textbf{Normal} \(\mathcal{N}(\mu, \sigma^2)\): \(M_X(t) = \exp\left(\mu t + \frac{1}{2}\sigma^2 t^2\right)\).
    \item \textbf{Bernoulli} (\(p\)): \(M_X(t) = 1 - p + p e^t\).
\end{itemize}

\subsection{Hoeffding's Lemma}
Let \(X\) be a random variable with \(\mathbb{E}[X] = \mu\) and \(\mathbb{P}(a \le X \le b) = 1\). Then for all \(t \in \mathbb{R}\):
\[
 \mathbb{E}[e^{t(X - \mu)}] \le \exp\left( \frac{t^2 (b-a)^2}{8} \right). 
\]

\emph{Proof.} Without loss of generality, assume \(\mu = 0\) (otherwise replace \(X\) with \(X-\mu\)). Since \(e^{tx}\) is convex, the graph of \(e^{tx}\) on \([a,b]\) lies below the secant line connecting \((a, e^{ta})\) and \((b, e^{tb})\). Thus,
\[
 e^{tx} \le \frac{b-x}{b-a}e^{ta} + \frac{x-a}{b-a}e^{tb}. 
\]
Taking expectations and using \(\mathbb{E}[X]=0\),
\[
 \mathbb{E}[e^{tX}] \le \frac{b}{b-a}e^{ta} - \frac{a}{b-a}e^{tb}. 
\]
Let \(p = -a/(b-a) \in [0,1]\) and \(u = t(b-a)\). Then \(b/(b-a) = 1-p\). The right-hand side is \((1-p)e^{-pu} + p e^{(1-p)u}\). Taking the logarithm and applying Taylor's theorem yields \(\log(\mathbb{E}[e^{tX}]) \le \frac{1}{8}t^2(b-a)^2\). \(\square\)

\hrulefill

\section{Sub-Gaussian and Sub-Exponential Variables}

\subsection{Sub-Gaussian Random Variables}
A random variable \(X\) with \(\mathbb{E}[X]=0\) is \textbf{sub-Gaussian} with variance proxy \(\sigma^2\) if:
\[
 \mathbb{E}[e^{tX}] \le e^{\frac{\sigma^2 t^2}{2}} \quad \forall t \in \mathbb{R}. 
\]
\textbf{Equivalent characterizations:}
\begin{enumerate}
    \item \textbf{Tail bound:} \(\mathbb{P}(|X| \ge t) \le 2e^{-t^2/(2\sigma^2)}\).
    \item \textbf{Moment growth:} \((\mathbb{E}[|X|^p])^{1/p} \le C\sigma \sqrt{p}\) for \(p \ge 1\).
    \item \textbf{Sub-Gaussian norm:} \(\|X\|_{\psi_2} = \sup_{p \ge 1} p^{-1/2} (\mathbb{E}[|X|^p])^{1/p} < \infty\).
\end{enumerate}

\emph{Example.} Any bounded random variable \(X \in [a,b]\) with \(\mathbb{E}[X]=0\) is sub-Gaussian with \(\sigma = (b-a)/2\) (by Hoeffding's Lemma).

\subsection{Sub-Exponential Random Variables}
A random variable \(X\) with \(\mathbb{E}[X]=0\) is \textbf{sub-exponential} with parameters \((\nu, \alpha)\) if:
\[
 \mathbb{E}[e^{tX}] \le e^{\frac{\nu^2 t^2}{2}} \quad \text{for } |t| \le \frac{1}{\alpha}. 
\]
\textbf{Equivalent characterizations:}
\begin{enumerate}
    \item \textbf{Tail bound:} \(\mathbb{P}(|X| \ge t) \le 2e^{-t/\alpha}\) for \(t \ge \nu^2/\alpha\).
    \item \textbf{Sub-exponential norm:} \(\|X\|_{\psi_1} = \sup_{p \ge 1} p^{-1} (\mathbb{E}[|X|^p])^{1/p} < \infty\).
\end{enumerate}

\emph{Example.} If \(Z \sim \mathcal{N}(0,1)\), then \(X = Z^2 - 1\) is sub-exponential.

\hrulefill

\section{Concentration Inequalities for Sums}

\subsection{Hoeffding's Inequality}
Let \(X_1, \dots, X_n\) be independent random variables such that \(a_i \le X_i \le b_i\) almost surely. Let \(\bar{X} = \frac{1}{n} \sum X_i\) and \(\mu = \mathbb{E}[\bar{X}]\). Then for any \(t > 0\):
\[
 \mathbb{P}(\bar{X} - \mu \ge t) \le \exp\left( - \frac{2n^2 t^2}{\sum_{i=1}^n (b_i - a_i)^2} \right). 
\]

\emph{Proof.} Apply the Chernoff bound to \(Z = \sum (X_i - \mathbb{E}[X_i])\). By independence and Hoeffding's Lemma,
\[
 \mathbb{E}[e^{sZ}] = \prod_{i=1}^n \mathbb{E}[e^{s(X_i - \mathbb{E}[X_i])}] \le \prod_{i=1}^n e^{\frac{s^2(b_i-a_i)^2}{8}} = e^{\frac{s^2}{8}\sum (b_i-a_i)^2}. 
\]
Setting \(s = 4nt / \sum (b_i-a_i)^2\) and applying the Chernoff bound yields the result. \(\square\)

\subsection{Chernoff Bounds for Independent Bernoulli Variables}
Let \(X = \sum_{i=1}^n X_i\) where \(X_i \sim \Bernoulli(p_i)\) independently, and let \(\mu = \mathbb{E}[X] = \sum p_i\). Then for any \(\delta > 0\):
\begin{enumerate}
    \item \textbf{Upper tail:} \(\mathbb{P}(X \ge (1+\delta)\mu) \le \exp\left( -\frac{\delta^2 \mu}{2+\delta} \right)\).
    \item \textbf{Lower tail:} For \(0 < \delta < 1\), \(\mathbb{P}(X \le (1-\delta)\mu) \le \exp\left( -\frac{\delta^2 \mu}{2} \right)\).
\end{enumerate}

\subsection{Bernstein's Inequality}
Let \(X_1, \dots, X_n\) be independent random variables with \(\mathbb{E}[X_i] = 0\), \(|X_i| \le M\) almost surely, and \(\Var(X_i) = \sigma_i^2\). Let \(\sigma^2 = \frac{1}{n}\sum \sigma_i^2\). Then for any \(t > 0\):
\[
 \mathbb{P}\left( \frac{1}{n}\sum X_i \ge t \right) \le \exp\left( - \frac{n t^2}{2\sigma^2 + 2Mt/3} \right). 
\]

\subsection{McDiarmid's Inequality}
Let \(f: \mathcal{X}^n \to \mathbb{R}\) satisfy the bounded difference property: for each \(i\), there exists \(c_i > 0\) such that \(|f(x) - f(x')| \le c_i\) whenever \(x\) and \(x'\) differ only in the \(i\)-th coordinate. Then for any \(t > 0\):
\[
 \mathbb{P}(f(X) - \mathbb{E}[f(X)] \ge t) \le \exp\left( - \frac{2t^2}{\sum_{i=1}^n c_i^2} \right). 
\]

\hrulefill

\section{Matrix Concentration}

\subsection{Matrix Chernoff Bound}
Let \(\mathbf{X}_1, \dots, \mathbf{X}_n\) be independent, symmetric, positive semi-definite \(d \times d\) matrices with \(\mathbb{E}[\mathbf{X}_i] = \mathbf{0}\) and \(\lambda_{\max}(\mathbf{X}_i) \le R\) almost surely. Let \(\mu_{\max} = \lambda_{\max}\left(\sum_{i=1}^n \mathbb{E}[\mathbf{X}_i^2]\right)\). Then for any \(\delta \in (0,1)\):
\[
 \mathbb{P}\left( \lambda_{\max}\left( \sum_{i=1}^n \mathbf{X}_i \right) \ge (1+\delta)\mu_{\max} \right) \le d \cdot \exp\left( -\frac{\delta^2 \mu_{\max}}{3R} \right). 
\]

\hrulefill

\section{Application: Discrepancy and the Probabilistic Method}

Let \(V\) be a set of \(n\) elements and \(\mathcal{F}\) a family of subsets. A coloring \(\chi: V \to \{-1,1\}\) is chosen uniformly at random. For a set \(S \in \mathcal{F}\), the discrepancy is \(d(S) = \sum_{v \in S} \chi(v)\). 

Since \(\chi(v)\) are independent Rademacher variables, \(d(S)\) is sub-Gaussian:
\[
 \mathbb{P}(|d(S)| \ge t) \le 2\exp\left( -\frac{t^2}{2|S|} \right) \le 2\exp\left( -\frac{t^2}{2n} \right). 
\]
By the union bound over \(m = |\mathcal{F}|\) sets:
\[
 \mathbb{P}\left( \max_{S \in \mathcal{F}} |d(S)| \ge t \right) \le 2m \exp\left( -\frac{t^2}{2n} \right). 
\]
Setting the right-hand side to \(1/2\) (to prove existence via the probabilistic method) yields \(t = O(\sqrt{n \log m})\). This establishes that any set system can be colored with discrepancy at most \(O(\sqrt{n \log m})\), a foundational result in discrepancy theory.

\end{document}
